% !TEX root = ../main.tex
\begin{abstractind}
  \justifying
  Peningkatan mobilitas kendaraan roda dua di lingkungan perguruan tinggi meningkatkan risiko keselamatan, khususnya terkait ketidakpatuhan penggunaan helm. Pengawasan manual konvensional dinilai kurang terukur dan belum mampu menyediakan wawasan spasiotemporal yang presisi untuk evaluasi kebijakan. Penelitian ini bertujuan merancang dan mengevaluasi arsitektur sistem pengawasan cerdas berbasis aliran video waktu nyata (\textit{real-time video streaming}) secara ujung ke ujung (\textit{end-to-end}). Arsitektur yang diusulkan mengintegrasikan model YOLOv8 untuk ekstraksi fitur spasial, ByteTrack untuk pelacakan multiobjek, serta filter \textit{Intersection over Area} (IoA) dan konfirmasi temporal berbasis \textit{N-Frame} untuk mereduksi deteksi palsu. Aliran peristiwa kemudian didistribusikan melalui Apache Kafka dan Apache Spark Structured Streaming ke \textit{data mart} PostgreSQL sebelum divisualisasikan pada dasbor interaktif Grafana. Hasil evaluasi eksperimental menunjukkan bahwa model deteksi mencapai \textit{Mean Average Precision} (mAP@50) sebesar 91,4\%. Sistem pemrosesan tepi beroperasi stabil dengan \textit{throughput} 18,9 FPS dan latensi inferensi 53 milidetik, sementara pialang pesan mencatat tingkat keberhasilan pengiriman 100\% dengan median latensi 2,65 milidetik. Pengujian operasional menggunakan data lapangan selama 10 jam (07.00--17.00 WIB) memvalidasi keandalan \textit{pipeline} dalam menyaring 11.172 deteksi mentah menjadi 113 kejadian pelanggaran aktual. Analisis visual pada dasbor berhasil memetakan distribusi kepadatan bimodal secara empiris, dengan puncak anomali ekstrem terekam pada pukul 14.45--15.15 WIB yang berkorelasi dengan faktor cuaca dan dinamika pergantian kelas. Secara keseluruhan, infrastruktur terintegrasi ini mentransformasi metode pengawasan reaktif menjadi instrumen Sistem Pendukung Keputusan (\textit{Decision Support System}) yang proaktif untuk mengoptimalkan efisiensi perencanaan patroli keselamatan di lingkungan kampus.

\noindent
\textbf{Kata kunci:} ByteTrack, \textit{Data Mart}, Pelanggaran Helm, YOLOv8 % masukkan keyword
\end{abstractind}
